\begin{enumerate}
\item \textbf{Start-up message.} On the MSX's own screen, \texttt{OCEAN SOFTWARE 1989 / DIHALT STUDIO 2026 / V9990 found !} for about two seconds.
\item \textbf{DiHalt logo.} Shown for four seconds; it cannot be skipped.
\item \textbf{Cover.} The illustration of the arcade game. Press any key or a joystick button to continue.
\item \textbf{Title screen.} \texttt{CHASE H.Q.} with a blinking \texttt{PRESS GEAR TO PLAY}; the credits of the MSX version and the copyright take turns underneath. Press \textbf{GEAR} (\key{SPACE} or joystick button 1) to start.
\item \textbf{Briefing.} The \texttt{CHASE H.Q. MONITORING SYSTEM} screen: Nancy, at Chase H.Q., tells you by radio who the criminal is, and the monitor shows his car. Press \textbf{GEAR} to skip it.
\item \textbf{The chase.} After each arrest the briefing of the next stage begins. There are five stages; after the fifth comes the ending, and then the title screen.
\end{enumerate}

Between screens the picture goes black for a moment while the next one is prepared. When the game is over, \texttt{GAME OVER} appears and the title screen comes back.
